\chapter{\MakeUppercase{Kesimpulan dan Saran}}

Bab ini berisi kesimpulan yang diperoleh dari hasil penelitian serta saran yang dapat digunakan sebagai masukan untuk pengembangan lebih lanjut. Kesimpulan dan saran harus disusun berdasarkan hasil analisis dan pembahasan yang telah dipaparkan pada bab sebelumnya.

\section{Kesimpulan}

Kesimpulan merupakan pernyataan ringkas yang merangkum hasil utama penelitian. Isi kesimpulan harus secara langsung menjawab rumusan masalah dan menunjukkan tingkat pencapaian tujuan penelitian yang telah ditetapkan.

Kesimpulan tidak diperkenankan memuat data baru, teori baru, atau pembahasan yang belum dijelaskan pada bab sebelumnya. Setiap kesimpulan harus didasarkan pada hasil pengujian, analisis, dan pembahasan yang telah dilakukan.

Dalam menyusun kesimpulan, beberapa hal berikut perlu diperhatikan:

\begin{itemize}[topsep=0pt,itemsep=0pt,parsep=0pt,leftmargin=*]
\item Kesimpulan harus menjawab seluruh rumusan masalah penelitian.
\item Kesimpulan harus menunjukkan apakah tujuan penelitian telah tercapai.
\item Kesimpulan disusun berdasarkan fakta dan hasil yang diperoleh selama penelitian.
\item Kesimpulan ditulis secara ringkas, jelas, dan tidak mengulang seluruh isi pembahasan.
\item Apabila memungkinkan, sertakan hasil atau capaian utama yang paling signifikan dari penelitian.
\end{itemize}

Kesimpulan dapat dituliskan dalam bentuk paragraf maupun butir-butir pernyataan sesuai dengan kebutuhan dan karakteristik penelitian.

\section{Saran}

Saran berisi rekomendasi yang diberikan berdasarkan pengalaman dan temuan selama pelaksanaan penelitian. Saran dapat ditujukan untuk pengembangan sistem, penyempurnaan metode, perluasan ruang lingkup penelitian, maupun penelitian lanjutan.

Saran harus bersifat konstruktif, realistis, dan relevan dengan hasil penelitian yang telah dilakukan. Saran tidak digunakan untuk memperbaiki kelemahan penulisan laporan, melainkan untuk memberikan arah pengembangan pada penelitian atau implementasi berikutnya.

Beberapa bentuk saran yang dapat diberikan antara lain:

\begin{itemize}[topsep=0pt,itemsep=0pt,parsep=0pt,leftmargin=*]
\item Pengembangan fitur atau fungsi yang belum diimplementasikan.
\item Peningkatan performa, akurasi, efisiensi, atau keandalan sistem.
\item Penggunaan metode, algoritma, atau pendekatan lain yang berpotensi memberikan hasil lebih baik.
\item Pengujian pada kondisi, lingkungan, atau skala yang lebih luas.
\item Pengembangan penelitian lanjutan untuk mengatasi keterbatasan yang masih ada.
\end{itemize}

Saran yang diberikan sebaiknya berkaitan langsung dengan keterbatasan, kendala, atau peluang pengembangan yang ditemukan selama penelitian berlangsung.
